\documentclass[11pt,a4paper]{article}
\input{preamble}

\pagestyle{fancy}
\fancyhf{}
\fancyhead[L]{\small\color{lightbrown}\emoji{brain} RL 틱택토 이론해설서}
\fancyhead[R]{\small\color{lightbrown}v1.0.0.0}
\fancyfoot[C]{\thepage}
\renewcommand{\headrulewidth}{0.4pt}
\setlength{\headheight}{14pt}

\begin{document}

\begin{titlepage}
\centering
\vspace*{3cm}
{\fontsize{48pt}{50pt}\selectfont \emoji{brain}}
\vspace{1cm}

{\Huge\bfseries\color{mainblue} RL 틱택토}
\vspace{0.3cm}

{\LARGE\color{darkbrown} 이론해설서}
\vspace{1.5cm}

{\large Q-Learning 강화학습의 원리와 개념 안내}
\vspace{3cm}

\begin{tabular}{rl}
\textbf{문서 버전} & v1.0.0.0 \\
\textbf{시뮬레이터} & rl-tictactoe-v3.html (v3) \\
\textbf{작성일} & 2026-02-14 \\
\end{tabular}
\vfill
{\small\color{lightbrown} 청주교육대학교 — 컴퓨팅사고와 AI교육}
\end{titlepage}

\tableofcontents
\newpage

% ============================================================
\section{강화학습이란?}
% ============================================================

\subsection{세 가지 학습 방법}

머신러닝에는 크게 세 가지 학습 방법이 있습니다.

\begin{center}
\renewcommand{\arraystretch}{1.4}
\begin{tabular}{|l|l|l|l|}
\hline
\rowcolor{tableheader}
\textbf{학습 방법} & \textbf{핵심} & \textbf{비유} & \textbf{예시} \\
\hline
지도학습 & 정답이 있는 데이터로 학습 & 선생님이 정답을 알려줌 & 사과 가격 예측, 과일 분류 \\
\hline
비지도학습 & 정답 없이 패턴을 찾음 & 스스로 분류하기 & K-means 군집화 \\
\hline
\textbf{강화학습} & \textbf{시행착오로 보상을 최대화} & \textbf{게임하며 스스로 배움} & \textbf{틱택토, 바둑, 로봇 제어} \\
\hline
\end{tabular}
\end{center}

강화학습(Reinforcement Learning, RL)은 \textbf{에이전트(agent)}가 \textbf{환경(environment)}과 상호작용하며 \textbf{보상(reward)}을 최대화하는 방법을 스스로 찾아가는 학습 방식입니다.

\subsection{강화학습의 핵심 요소}

\begin{center}
\renewcommand{\arraystretch}{1.4}
\begin{tabular}{|l|l|p{6cm}|}
\hline
\rowcolor{tableheader}
\textbf{요소} & \textbf{틱택토 대응} & \textbf{설명} \\
\hline
에이전트 (Agent) & AI 플레이어 (X 또는 O) & 행동을 결정하는 학습 주체 \\
\hline
환경 (Environment) & 틱택토 보드 & 에이전트가 상호작용하는 세계 \\
\hline
상태 (State) & 현재 보드 배치 & 환경의 현재 모습 \\
\hline
행동 (Action) & 빈 칸에 돌 놓기 & 에이전트가 취할 수 있는 선택 \\
\hline
보상 (Reward) & 승리(+1), 패배($-$1), 무승부(+0.3) & 행동의 결과에 대한 피드백 \\
\hline
정책 (Policy) & Q-테이블 기반 선택 & 상태에서 어떤 행동을 할지의 전략 \\
\hline
\end{tabular}
\end{center}

\subsection{강화학습의 학습 흐름}

에이전트가 현재 보드(상태)를 보고 $\to$ 빈 칸 중 하나를 골라 돌을 놓고(행동) $\to$ 결과가 나오면(보상) $\to$ ``이 상황에서 이 칸에 두니까 이겼구나/졌구나'' 경험을 기록하고 $\to$ 다음에 비슷한 상황이 오면 더 나은 선택을 합니다. 이 과정을 수천\textasciitilde수만 번 반복하면 에이전트는 점점 더 좋은 전략을 학습합니다.

% ============================================================
\section{Q-Learning 알고리즘}
% ============================================================

\subsection{Q-값이란?}

Q-값(Q-value)은 \textbf{``이 상태에서 이 행동을 하면 앞으로 얼마나 좋을까?''}를 숫자로 나타낸 것입니다.

\begin{tipbox}[\emoji{light-bulb} 쉬운 비유]
음식점 평점과 비슷해요. ``이 가게에서 파스타를 시키면 4.5점, 피자를 시키면 3.8점'' — 그러면 파스타를 주문하겠죠? Q-값도 ``이 보드에서 중앙에 두면 0.42점, 모서리에 두면 0.18점'' 이런 식이에요.
\end{tipbox}

\begin{center}
\renewcommand{\arraystretch}{1.4}
\begin{tabular}{|c|l|l|}
\hline
\rowcolor{tableheader}
\textbf{기호} & \textbf{읽는 법} & \textbf{의미} \\
\hline
$Q(s, a)$ & ``큐 에스 에이'' & 상태 $s$에서 행동 $a$의 가치 \\
\hline
$s$ & state & 현재 보드 배치 \\
\hline
$a$ & action & 돌을 놓을 위치 (0\textasciitilde8번 칸) \\
\hline
\end{tabular}
\end{center}

\subsection{Q-테이블}

Q-테이블은 \textbf{모든 (상태, 행동) 쌍에 대한 Q-값을 저장하는 표}입니다. 시뮬레이터에서 \textbf{\emoji{books} Q-상태} 지표는 이 테이블에 몇 개의 상태가 기록되었는지를 보여줍니다.

\subsection{Q-값 갱신 규칙}

게임이 끝나면 에이전트는 자신이 둔 수들을 \textbf{마지막 수부터 거꾸로} 돌아보며 Q-값을 업데이트합니다.

\begin{infobox}[갱신 공식]
\centering\Large
$Q(s, a) \leftarrow Q(s, a) + \alpha \times (\text{보상} - Q(s, a))$
\end{infobox}

\begin{center}
\renewcommand{\arraystretch}{1.4}
\begin{tabular}{|c|l|l|c|}
\hline
\rowcolor{tableheader}
\textbf{기호} & \textbf{이름} & \textbf{의미} & \textbf{기본값} \\
\hline
$\alpha$ & 학습률 (alpha) & 새 경험을 얼마나 반영할지 & 0.30 \\
\hline
보상 & reward & 해당 수의 결과에 대한 평가 & — \\
\hline
\end{tabular}
\end{center}

\begin{tipbox}[\emoji{light-bulb} 쉬운 비유]
``기존 점수에 새 경험을 30\%만 반영한다''는 뜻이에요. $\alpha=0.3$이면 기존 생각 70\% + 새 경험 30\%로 평가를 업데이트해요.
\end{tipbox}

\subsection{역방향 갱신과 할인}

시뮬레이터는 에피소드가 끝나면 \textbf{마지막 수 $\to$ 첫 수} 순서로 거슬러 올라가며 갱신합니다. 이때 \textbf{할인율($\gamma$)}이 적용됩니다.

$\gamma=0.9$이면 한 수 뒤의 보상은 90\%로 줄어들고, 두 수 뒤는 81\%로 줄어듭니다. \textbf{가까운 미래의 보상이 더 중요}하다는 뜻입니다.

\subsection{갱신 예시}

X가 승리($+1.0$)한 에피소드에서 X가 3수를 뒀다고 가정합니다.

\begin{codebox}
3수째 (마지막): reward $= +1.0$

\quad $Q(s_3, a_3) \leftarrow Q(s_3, a_3) + 0.3 \times (1.0 - Q(s_3, a_3))$

2수째: reward $= -0.01 + 0.9 \times 1.0 = +0.89$

\quad $Q(s_2, a_2) \leftarrow Q(s_2, a_2) + 0.3 \times (0.89 - Q(s_2, a_2))$

1수째 (첫 수): reward $= -0.01 + 0.9 \times 0.89 = +0.80$

\quad $Q(s_1, a_1) \leftarrow Q(s_1, a_1) + 0.3 \times (0.80 - Q(s_1, a_1))$
\end{codebox}

승리한 경우 Q-값이 양의 방향으로 올라가고, 패배하면 음의 방향으로 내려갑니다.

% ============================================================
\section{탐험과 활용의 균형}
% ============================================================

\subsection{$\varepsilon$-greedy 전략}

에이전트는 매 수마다 두 가지 중 하나를 선택합니다.

\begin{center}
\renewcommand{\arraystretch}{1.4}
\begin{tabular}{|c|l|l|}
\hline
\rowcolor{tableheader}
\textbf{확률} & \textbf{선택 방식} & \textbf{이름} \\
\hline
$\varepsilon$ & 빈 칸 중 \textbf{무작위} 선택 & \emoji{game-die} 탐험 (exploration) \\
\hline
$1 - \varepsilon$ & Q-값이 \textbf{가장 높은} 칸 선택 & \emoji{bullseye} 활용 (exploitation) \\
\hline
\end{tabular}
\end{center}

\begin{tipbox}[\emoji{light-bulb} 쉬운 비유]
매일 점심을 먹을 때, 가끔은 새로운 가게를 시도해보고(탐험), 대부분은 가장 맛있다고 알려진 가게에 가는(활용) 것과 같아요.
\end{tipbox}

\subsection{왜 탐험이 필요한가?}

처음부터 활용만 하면 \textbf{``처음 우연히 이긴 수만 계속 두는''} 문제가 생깁니다. 아직 시도하지 않은 더 좋은 수가 있을 수 있기 때문이죠. 탐험은 다양한 상태를 경험하게 해서 Q-테이블을 풍부하게 채우는 역할을 합니다.

\subsection{$\varepsilon$의 감소}

학습 초반에는 탐험을 많이 하고, 점차 줄여나갑니다. 초기 $\varepsilon = 100\%$ (거의 모든 수를 랜덤으로) $\to$ 매 배치마다 $\varepsilon \leftarrow \varepsilon \times 0.9995$ $\to$ 중기 $\varepsilon \approx 30\%$ $\to$ 후기 $\varepsilon \approx 1\%$ (거의 항상 최적수 선택).

\subsection{감소율의 의미}

\begin{center}
\renewcommand{\arraystretch}{1.4}
\begin{tabular}{|c|c|l|}
\hline
\rowcolor{tableheader}
\textbf{$\varepsilon$ 감소율} & \textbf{1\%까지 도달} & \textbf{특징} \\
\hline
0.990 & $\sim$460판 & 매우 빠른 전환 — 적은 탐험 \\
\hline
0.995 & $\sim$920판 & 빠른 전환 \\
\hline
0.9995 (기본) & $\sim$9,200판 & 적당한 탐험 기간 \\
\hline
0.9999 & $\sim$46,000판 & 긴 탐험 — 더 많은 상태 발견 \\
\hline
\end{tabular}
\end{center}

% ============================================================
\section{보상 설계}
% ============================================================

\subsection{보상 체계}

\begin{center}
\renewcommand{\arraystretch}{1.4}
\begin{tabular}{|l|c|p{7cm}|}
\hline
\rowcolor{tableheader}
\textbf{결과} & \textbf{보상값} & \textbf{설계 이유} \\
\hline
승리 & $+1.0$ & 가장 바람직한 결과 \\
\hline
패배 & $-1.0$ & 가장 피해야 할 결과 \\
\hline
무승부 & $+0.3$ & 패배보다 낫지만 승리만큼은 아님 \\
\hline
매 수 & $-0.01$ & 빠른 승리를 유도하는 단계 패널티 \\
\hline
\end{tabular}
\end{center}

\subsection{무승부 보상이 양수인 이유}

무승부 보상을 0으로 하면, 에이전트는 ``이길 수 없는 상황에서 어떻게 해야 할지'' 학습하기 어렵습니다. $+0.3$을 부여하면 패배($-1.0$)를 피하고 무승부($+0.3$)를 선택하는 \textbf{방어적 플레이}를 학습합니다.

\subsection{단계 패널티의 역할}

매 수마다 $-0.01$의 패널티가 있으면, AI는 \textbf{같은 결과라도 더 빨리} 달성하려 합니다. 3수 만에 이기면 패널티가 $-0.03$이지만, 5수 만에 이기면 $-0.05$입니다.

% ============================================================
\section{대칭성 정규화}
% ============================================================

\subsection{대칭 변환이란?}

틱택토 보드는 \textbf{8가지 대칭 변환}(원본, 90°, 180°, 270° 회전 $\times$ 좌우 반사)으로 같은 모양을 만들 수 있습니다. 위 8가지 보드는 \textbf{전략적으로 완전히 동일}합니다.

\subsection{정규화의 효과}

\begin{center}
\renewcommand{\arraystretch}{1.4}
\begin{tabular}{|l|c|c|c|}
\hline
\rowcolor{tableheader}
\textbf{항목} & \textbf{정규화 없음} & \textbf{정규화 적용} & \textbf{효율} \\
\hline
빈 보드의 첫 수 선택지 & 9개 & 3개 (중앙, 모서리, 변) & 3배 \\
\hline
전체 Q-상태 수 & $\sim$5,000개 & $\sim$765개 & $\sim$7배 \\
\hline
같은 성능 도달 시간 & 기준 & 약 $1/7\sim1/8$ & 7\textasciitilde8배 \\
\hline
\end{tabular}
\end{center}

\subsection{정규화 작동 방식}

현재 보드를 8가지로 변환하고, 사전식 순서가 가장 작은 것을 ``정규 상태''로 선택합니다. 정규 좌표계에서 Q-값을 조회하고 행동을 선택한 뒤, 선택한 행동을 원본 좌표로 역변환하여 실행합니다. 충분히 학습하면 Q-상태 수가 약 765개 부근에서 포화됩니다.

% ============================================================
\section{학습 과정의 이해}
% ============================================================

\subsection{학습 5단계}

\begin{center}
\renewcommand{\arraystretch}{1.4}
\small
\begin{tabular}{|l|c|p{3.5cm}|p{4cm}|}
\hline
\rowcolor{tableheader}
\textbf{단계} & \textbf{에피소드} & \textbf{무엇이 일어나는가} & \textbf{관찰 포인트} \\
\hline
\emoji{seedling} 초기 & 0\textasciitilde499 & Q-테이블 비어있어 거의 랜덤 & X/O 승률 비슷, $\varepsilon\approx100\%$ \\
\hline
\emoji{magnifying-glass-tilted-left} 탐색 & 500\textasciitilde4,999 & 기본 패턴 발견, 탐험 감소 & X 승률 서서히 상승 \\
\hline
\emoji{books} 학습 & 5,000\textasciitilde19,999 & 전략 구체화, 무승부율 상승 & $\varepsilon<20\%$, 강도 50+ \\
\hline
\emoji{bullseye} 응용 & 20,000\textasciitilde49,999 & 안정적 전략, 미세 조정 & 무승부율 80\%+, 강도 80+ \\
\hline
\emoji{trophy} 숙련 & 50,000+ & 이론적 최적에 근접 & 무승부율 $\sim$100\% \\
\hline
\end{tabular}
\end{center}

\subsection{승률 추이의 해석}

학습이 진행되면 특징적인 패턴이 나타납니다. 초기에는 X 승률 $\sim$40\%, O 승률 $\sim$30\%, 무승부 $\sim$30\%이던 것이, 중기에는 무승부 $\sim$65\%로 올라가고, 학습 완료 시 무승부 $\sim$100\%에 수렴합니다. \textbf{양쪽 AI 모두 최적 학습을 마치면 무승부율이 100\%에 수렴합니다.}

\subsection{첫수 선호도의 변화}

\begin{center}
\renewcommand{\arraystretch}{1.4}
\begin{tabular}{|l|l|l|}
\hline
\rowcolor{tableheader}
\textbf{학습 단계} & \textbf{선호 순서} & \textbf{이유} \\
\hline
초기 & 랜덤 (모두 $Q\approx0$) & 아직 경험 부족 \\
\hline
중기 & 중앙 > 모서리 $\approx$ 변 & 중앙의 이점 발견 \\
\hline
완료 & \emoji{bullseye} 중앙 >> 모서리 >> 변 & 최적 전략 확립 \\
\hline
\end{tabular}
\end{center}

% ============================================================
\section{틱택토의 게임 이론}
% ============================================================

\subsection{완전 정보 게임}

틱택토는 \textbf{완전 정보(perfect information) 제로섬(zero-sum) 유한(finite) 게임}입니다. 모든 플레이어가 보드 전체를 볼 수 있고(완전 정보), 한쪽이 이기면 다른 쪽은 반드시 지며(제로섬), 게임이 반드시 끝납니다(유한, 최대 9수).

\subsection{게임 트리}

\begin{center}
\renewcommand{\arraystretch}{1.4}
\begin{tabular}{|l|c|}
\hline
\rowcolor{tableheader}
\textbf{항목} & \textbf{값} \\
\hline
가능한 게임 진행 수 (대칭 무시) & 255,168가지 \\
\hline
서로 다른 보드 상태 & 5,478개 \\
\hline
대칭 정규화 후 상태 & $\sim$765개 \\
\hline
이론적 결론 & \textbf{양쪽 최적 시 무승부} \\
\hline
\end{tabular}
\end{center}

\subsection{최적 전략}

선공(X)의 최적 첫수: \emoji{bullseye} \textbf{중앙}(4번 칸)이 유일한 최적 선택입니다.

\begin{center}
\renewcommand{\arraystretch}{1.4}
\begin{tabular}{|l|c|c|}
\hline
\rowcolor{tableheader}
\textbf{위치 유형} & \textbf{관여하는 승리선 수} & \textbf{전략적 가치} \\
\hline
\emoji{bullseye} 중앙 (4번) & 4개 & 최고 \\
\hline
모서리 (0, 2, 6, 8번) & 3개 & 높음 \\
\hline
변 (1, 3, 5, 7번) & 2개 & 낮음 \\
\hline
\end{tabular}
\end{center}

\subsection{Q-Learning과 미니맥스의 관계}

\begin{center}
\renewcommand{\arraystretch}{1.4}
\begin{tabular}{|l|c|c|}
\hline
\rowcolor{tableheader}
\textbf{비교} & \textbf{Q-Learning} & \textbf{미니맥스} \\
\hline
접근 & 경험으로 학습 & 전체 탐색 \\
\hline
사전 지식 & 불필요 & 게임 규칙 필요 \\
\hline
최적성 & 충분한 학습 시 근사 & 항상 최적 \\
\hline
확장성 & 큰 게임에도 적용 가능 & 큰 게임은 계산 불가 \\
\hline
교육적 가치 & 학습 과정 관찰 가능 & 결과만 확인 가능 \\
\hline
\end{tabular}
\end{center}

% ============================================================
\section{파라미터가 학습에 미치는 영향}
% ============================================================

\subsection{학습률 $\alpha$ (Alpha)}

$\alpha$는 \textbf{새로운 경험을 얼마나 강하게 반영할지} 결정합니다. $\alpha$가 작으면(0.01): 오래된 경험 99\% + 새 경험 1\% $\to$ 안정적이지만 느린 학습. $\alpha$가 크면(0.9): 오래된 경험 10\% + 새 경험 90\% $\to$ 빠르지만 불안정한 학습.

\subsection{할인율 $\gamma$ (Gamma)}

$\gamma$는 \textbf{미래의 보상을 현재에 얼마나 반영할지} 결정합니다.

\begin{center}
\renewcommand{\arraystretch}{1.4}
\begin{tabular}{|c|c|l|}
\hline
\rowcolor{tableheader}
\textbf{$\gamma$ 값} & \textbf{n수 뒤 반영률} & \textbf{전략 성향} \\
\hline
0.1 & 1수: 10\%, 2수: 1\% & 눈앞의 이익만 추구 \\
\hline
0.5 & 1수: 50\%, 2수: 25\% & 1\textasciitilde2수 앞을 봄 \\
\hline
0.9 & 1수: 90\%, 2수: 81\% & 여러 수 앞을 봄 \\
\hline
1.0 & 1수: 100\%, 2수: 100\% & 모든 미래를 동등하게 평가 \\
\hline
\end{tabular}
\end{center}

\subsection{$\varepsilon$ 감소율 (Epsilon Decay)}

감소율은 \textbf{탐험에서 활용으로 전환하는 속도}를 결정합니다. 0.990은 매우 짧은 탐험($\sim$460판), 0.9995(기본)는 적당($\sim$9,200판), 0.9999는 길음($\sim$46,000판).

\subsection{파라미터 상호작용}

\begin{center}
\renewcommand{\arraystretch}{1.4}
\begin{tabular}{|l|p{8cm}|}
\hline
\rowcolor{tableheader}
\textbf{조합} & \textbf{결과} \\
\hline
높은 $\alpha$ + 낮은 $\varepsilon$ 감소 & 빠르지만 불안정 — 그래프가 심하게 요동 \\
\hline
낮은 $\alpha$ + 높은 $\varepsilon$ 감소 & 느리지만 안정 — 그래프가 천천히 올라감 \\
\hline
높은 $\gamma$ + 긴 탐험 & 장기 전략 학습에 충분한 시간 \\
\hline
낮은 $\gamma$ + 짧은 탐험 & 단기 전략으로 빠르게 수렴 \\
\hline
\end{tabular}
\end{center}

% ============================================================
\section{시뮬레이터와 이론의 연결}
% ============================================================

\subsection{개념-UI 매핑}

\begin{center}
\renewcommand{\arraystretch}{1.3}
\small
\begin{tabular}{|l|l|p{4.5cm}|}
\hline
\rowcolor{tableheader}
\textbf{이론 개념} & \textbf{시뮬레이터 UI} & \textbf{관찰 방법} \\
\hline
Q-값 & 보드 위 숫자, 히트맵 색상 & 사람 대전에서 빈 칸의 숫자 확인 \\
\hline
Q-테이블 크기 & \emoji{books} Q-상태 수 & 대시보드 학습 통계 \\
\hline
$\varepsilon$ (탐험률) & \emoji{game-die} $\varepsilon$ 카드, 그래프 & 100\%에서 1\%로 감소 관찰 \\
\hline
$\alpha$ (학습률) & \emoji{gear} 슬라이더 & 값을 바꾸고 승률 그래프 변화 관찰 \\
\hline
보상 & 학습 통계의 Step 패널티 & $-0.01$ 표시 확인 \\
\hline
정책 & Q-값 히트맵의 \emoji{star} 최선 & AI가 선택하는 최적수 \\
\hline
대칭 정규화 & Q-상태 수 $\sim$765 & 비정규화 시 $\sim$5,000과 비교 \\
\hline
학습 수렴 & 무승부율 그래프 & 100\%에 수렴하는 과정 \\
\hline
\end{tabular}
\end{center}

% ============================================================
\section{더 알아보기}
% ============================================================

\subsection{Q-Learning의 한계}

\begin{center}
\renewcommand{\arraystretch}{1.4}
\begin{tabular}{|l|p{5cm}|l|}
\hline
\rowcolor{tableheader}
\textbf{한계} & \textbf{설명} & \textbf{대안} \\
\hline
상태 공간 폭발 & 바둑 같은 큰 게임은 Q-테이블 저장 불가 & Deep Q-Network \\
\hline
연속 상태 & 로봇 제어 등 연속적 상태에서 테이블 불가 & 함수 근사 방법 \\
\hline
단일 에이전트 & 여러 에이전트 협력/경쟁은 다름 & 멀티 에이전트 RL \\
\hline
\end{tabular}
\end{center}

\subsection{강화학습의 실제 응용}

\begin{center}
\renewcommand{\arraystretch}{1.4}
\begin{tabular}{|l|p{9cm}|}
\hline
\rowcolor{tableheader}
\textbf{분야} & \textbf{사례} \\
\hline
게임 & AlphaGo(바둑), OpenAI Five(Dota2), AlphaStar(스타크래프트) \\
\hline
로봇 & 보행 학습, 물건 잡기, 자율 주행 \\
\hline
추천 & 뉴스/영상 추천 알고리즘 \\
\hline
의료 & 치료 순서 최적화 \\
\hline
에너지 & 데이터센터 냉각 최적화 (DeepMind) \\
\hline
\end{tabular}
\end{center}

\subsection{용어 사전}

\begin{center}
\renewcommand{\arraystretch}{1.3}
\small
\begin{tabular}{|l|l|p{6.5cm}|}
\hline
\rowcolor{tableheader}
\textbf{한국어} & \textbf{영어} & \textbf{의미} \\
\hline
에이전트 & Agent & 학습하는 AI 주체 \\
\hline
환경 & Environment & 에이전트가 상호작용하는 세계 \\
\hline
상태 & State & 환경의 현재 모습 \\
\hline
행동 & Action & 에이전트의 선택 \\
\hline
보상 & Reward & 행동 결과에 대한 피드백 \\
\hline
에피소드 & Episode & 한 판의 게임 (시작$\sim$끝) \\
\hline
Q-값 & Q-value & 상태-행동 쌍의 기대 가치 \\
\hline
Q-테이블 & Q-table & 모든 Q-값을 저장한 표 \\
\hline
정책 & Policy & 상태에서 행동을 결정하는 규칙 \\
\hline
탐험 & Exploration & 새로운 행동 시도 \\
\hline
활용 & Exploitation & 최선으로 알려진 행동 선택 \\
\hline
할인율 & Discount factor & 미래 보상의 현재 가치 비율 \\
\hline
수렴 & Convergence & Q-값이 최적에 가까워지는 것 \\
\hline
\end{tabular}
\end{center}

% ============================================================
\newpage
\section*{참고자료}
% ============================================================

\begin{center}
\renewcommand{\arraystretch}{1.4}
\begin{tabular}{|c|l|l|l|}
\hline
\rowcolor{tableheader}
\textbf{\#} & \textbf{자료명} & \textbf{유형} & \textbf{비고} \\
\hline
1 & RL틱택토\_개발일지\_v1.0.0.0 & 프로젝트 문서 & 알고리즘 구현 상세 \\
\hline
2 & RL틱택토\_사용설명서\_v1.0.0.0 & 프로젝트 문서 & UI 기능 상세 \\
\hline
3 & RL틱택토\_퀵가이드\_v1.0.0.0 & 프로젝트 문서 & 빠른 참조 \\
\hline
4 & Sutton \& Barto (2018) & 교재 & 강화학습 표준 교재 \\
\hline
5 & Watkins (1989) & 논문 & Q-Learning 원논문 \\
\hline
\end{tabular}
\end{center}

\section*{생성/수정 이력}

\begin{center}
\renewcommand{\arraystretch}{1.3}
\begin{tabular}{|l|l|l|l|p{5.5cm}|l|}
\hline
\rowcolor{tableheader}
\textbf{버전} & \textbf{날짜} & \textbf{시간} & \textbf{변경 수준} & \textbf{변경 내용} & \textbf{작성자} \\
\hline
v1.0.0.0 & 2026-02-14 & 21:30 & 최초 & 초안 작성 — 10개 섹션 전체 & Claude \\
\hline
\end{tabular}
\end{center}

\vfill
\begin{center}
\small\color{lightbrown}
\emph{RL 틱택토 이론해설서 — Q-Learning 강화학습의 원리와 개념 안내}
\end{center}

\end{document}
